
\subsection{Evaluating}

\subsubsection{Allocation Morphisms}

We call morphisms in $\tempC$ and $\tinstC$ that are both variable preserving and reflecting as allocation morphisms.
We intent to capture the idea that allocations only introduce concrete locations and do not interfere with location variables.

\begin{definition}[Allocation Morphisms]
    A morphism $\rho : \sigma_1 \to \sigma_2$ is an allocation morphism if it satisfies 
        $\concL{\sigma_2} \pltimes{\anyL{\sigma_2}} \varL{\sigma_1} = \emptyset$ and
        $\varL{\sigma_2} \pltimes{\anyL{\sigma_2}} \varL{\sigma_1} = \emptyset$.
\end{definition}
\begin{example}
    The morphism $?_\va \xto{\pair{?_\vb}{\incl^{\va \inclto \va,\vb}}} \incl^{\vb \inclto \va,\vb}$ is an allocation morphism since it does not instantiate nor introduce any variables.
\end{example}

\begin{notation}
    An allocation morphism is represented using the arrow $\ato$.
\end{notation}

\begin{lemma}[Composition of Allocation Morphisms]\label{lem:comp-alloc}
    Allocation morphisms are closed under composition.
\end{lemma}
\begin{proof}
    Follows from \cref{lem:comp-vpres} and \cref{lem:comp-vrefl}.
\end{proof}

\begin{lemma}[Codomain of Allocation Morphisms]\label{lem:codomain-alloc}
    If $\rho : \sigma_1 \to \sigma_2$ is an allocation morphism, then 
    \[ \sigma_2 \iso \begin{tikzcd}
                        {\concL{\sigma_2} \mplus \varL{\sigma_1}} \\
                        {\concL{\sigma_2}} \ar[u, "\incl_1"']
                    \end{tikzcd}\]
\end{lemma}
\begin{proof} We have the canonical isomorphism
    \[ \sigma_2 \iso \begin{tikzcd}
                    {\concL{\sigma_2} \mplus \varL{\sigma_2}} \\
                    {\concL{\sigma_2}} \ar[u, "\incl_1"']
                 \end{tikzcd}\]
    and notice that $\varL{\sigma_1} \iso \varL{\sigma_2}$, 
    from $\varL\rho \comp \inv^{\varL\rho}\res|{\sigma_2} \comp \varL\rho\res|{\sigma_2} = \varL\rho \comp 1$ and
    $\varL{\sigma_2} \comp \varL\rho\res|{\sigma_2} \comp \inv^{\varL\rho}\res|{\sigma_2} = \varL{\sigma_2} \comp 1$.
\end{proof}
\begin{notation}
    If $\rho$ is an allocation morphism, let $\aiso\rho$ be the isomorphism given by \cref{lem:codomain-alloc}.
\end{notation}

\begin{lemma}[Local Independent Coproducts and Allocation Morphisms]\label{lem:lic-alloc}
    In a L.I.C. square where $\rho_2$ is an allocation morphism
    \begin{center}
    \begin{tikzcd}[sep=small]     
                                                     & {\sigma} \ar[ld, "{\rho_1}"'] \ar[-alloctip, rd, "{\rho_2}" black]  & \\
        {\sigma_1} \ar[rd, "{\ext^{\rho_1}\rho_2}"'] &                                                                     & {\sigma_2} \ar[ld, "{\ext^{\rho_2}\rho_1}"] \\
                                                     &              {\sigma_1 \liplus\sigma \sigma_2}                          &
    \end{tikzcd}
    \end{center}
    we have that $\ext^{\rho_1}\rho_2$ is an allocation morphism as well.
\end{lemma}
\begin{proof}
    Follows from \cref{lem:lic-vpres} and \cref{lem:lic-vrefl}.
\end{proof}

\begin{lemma}[Allocation Morphism Cancellation]\label{lem:alloc-cancl}
    In \begin{tikzcd}
        {\sigma_1} 
            \ar[rr, "\rho"] \ar[dr, "\rho_1"']
                && {\sigma_3}
        \\
            & {\sigma_2}
                \ar[ur, "\rho_2"']
    \end{tikzcd}
    if $\rho$ is an allocation morphism then $\rho_1$ is an allocation morphism iff $\rho_2$ is.
\end{lemma}
\begin{proof}
    For ($\to$), it follows from \cref{lem:vrefl-cancl,lem:vpres-vrefl-cancl}, and for ($\from$) it follows from \cref{lem:vpres-cancl,lem:vrefl-vpres-cancl}.
\end{proof}

\begin{lemma}[Allocation Morphisms and Epis]\label{lem:alloc-epi}
    Every $\sigma_1 \xato{\me_1} \sigma_2 \xsurjto{\me_2} \sigma_3$ has the following diagram in $\tempC$
    \begin{center}
    \begin{tikzcd}
        {\sigma_{\uncon{\me_2 \comp \me_1}}}
            \ar[r, dashed, "\exists"] \ar[d, "\compl{\uncon{\me_2 \comp \me_1}}"']
            & {\sigma_2 \ominus \me_1}
                \ar[d, "\compl{\me_1}"]
        \\
        {\sigma_3}
            & {\sigma_2}
                \ar[l, "\me_2"]
    \end{tikzcd}
    \end{center}
\end{lemma}
\begin{proof}
    Notice that if it holds at $\anyL\dash$ level it also holds at $\concL\dash$ from 
    \begin{center}
    \begin{tikzcd}
            & {\concL{\sigma_2}}
                \ar[dr, "\concL{\sigma_2}"]
        \\
        {\concL{\sigma_2} \ominus \concL{\me_1}}
            \ar[ur, "\compl{\concL{\me_1}}"]
            & {\anyL{\sigma_2} \ominus \anyL{\me_1}}
                \ar[l, "\inv^{\sigma_2}\res|{\me_1}"'] \ar[r, "\compl{\anyL{\me_1}}"] 
                & {\anyL{\sigma_2}}
                    \ar[d, "\anyL{\me_2}"'] 
        \\
        {\concL{\sigma_3} \ominus \concL{\uncon{\me_2 \comp \me_1}}}
            \ar[r, "\concL{\sigma_3}\res|{\uncon{\me_2 \comp \me_1}}"' yshift=-5pt]
            & {\anyL{\sigma_3} \ominus \anyL{\uncon{\me_2 \comp \me_1}}}
                \ar[u, dashed, "\exists \inv^{\anyL{\me_2}}\res|{\me_1}"]
                \ar[r, "\compl{\anyL{\uncon{\me_2 \comp \me_1}}}"' yshift=-5pt] 
                & {\anyL{\sigma_3}}
    \end{tikzcd}
    \end{center}

    We construct $\anyL{\sigma_3} \ominus \anyL{\uncon{\me_2 \comp \me_1}} \dashto \anyL{\sigma_2} \ominus \concL{\me_1}$, 
    using the canonical iso
    \[ (\anyL{\sigma_{\con{\me_2 \comp \me_1}}} \ominus \anyL{\con{\me_2 \comp \me_1}}) \oplus 
       (\anyL{\sigma_3} \ominus \anyL{\uncon{\me_2 \comp \me_1}}) \iso (\anyL{\sigma_2} \ominus \anyL{\me_1}) \oplus (\anyL{\sigma_3} \ominus \anyL{\me_2}) \]
    Observe that if there exists $\va, \va' : \vC$ s.t. $\anyL{\compl{\uncon{\me_2 \comp \me_1}}}(\va) = \compl{\anyL{\me_2}}(\va')$
    then there exists $\vb : \vC$, s.t. $\anyL{\uncon{\me_2}}\compl{\anyL{\con{\me_2}}}(\vb) = \va'$ since $\me_2$ is connected. 
    As a result, we have a $\vb : \vC$ that is not in the image of $\anyL{\sigma_2}$ and is reachable in $\me_2$.
    By \cref{def:reach-loc-hwalks}, we have a heap walk that leads to $\vb$ with root location variable in $\varL{\sigma_2}$ and since $\me_1$ is an allocation morphism
    by \cref{lem:vrefl-prop}, this variable comes from $\varL{\sigma_1}$. 
    As a result, $\vb$ must be reachable in $\me_2 \comp \me_1$ which is a contradiction as it maps to the unreachable segment by $\anyL{\uncon{\me_2}}\compl{\anyL{\con{\me_2}}}$.
\end{proof}

\begin{lemma}[Disconnectedness Preservation and L.I.G.U.]\label{lem:dcon-ligu}
    A diagram of the form
    \begin{tikzcd}
        & {\sigma_1} 
            \ar[dl, twoheadrightarrow, "\me_1'"'] \ar[dr, -alloctip, "\me_1"]
    \\
    {\sigma_2'}
        \ar[dr, -rmonotip, "\me_2'"']
            && {\sigma_2}
                    \ar[dl, twoheadrightarrow, "\me_2"]
    \\
        & {\sigma_3}
    \end{tikzcd}
    s.t. $\me_2$ preserves the disconnectivity of $\me_1$ implies that it is a L.I.G.U.
\end{lemma}
\begin{proof}
    We first show that $\me_2'$ is variable reflecting.
    From \cref{lem:alloc-epi,lem:vrefl-prop} we have
    \begin{center}
        \begin{tikzcd}
            {\anyL{\sigma_3} \ominus \anyL{\me_2'}} 
                \ar[r, dashed] 
                    & {\anyL{\sigma_2} \ominus \anyL{\me_1}}
                        \ar[d, "\inv^{\concL{\sigma_2}}\res|{\me_1}"]
            \\
            {\concL{\sigma_3} \ominus \concL{\me_2'}} 
                \ar[r, dashed] \ar[u, "\concL{\sigma_3}\res|{\me_2'}"]
                    & {\concL{\sigma_2} \ominus \concL{\me_1}} 
    \end{tikzcd}
    \end{center}
    which gives us a morphism $\anyL{\sigma_3} \ominus \anyL{\me_2'} \to \concL{\sigma_2} \ominus \concL{\me_1}$ and notice that
    post-composing this with $\concL{\me_2} \comp \compl{\concL{\me_1}}$ must reside in $\concL{\sigma_3} \ominus \concL{\me_2'}$ 
    for otherwise the square of $\me_2'$ would not commute.

    By \cref{lem:lic-vpres}, we need to show that both $\concL\dash$ and $\anyL\dash$ images are L.I.C. 
    Notice that by \cref{lem:dcon-pres-prop-equiv}, its $\anyL\dash$-image is a pullback, 
    and by \cref{lem:pullback-vpres} we have that the $\tinstU$-image is a pullback in $\tempC$,
    and so the $\concL\dash$-image is a pullback as well. 

    So it enough to show that the $\concL\dash$ and $\anyL\dash$ images of cospan $\sigma_2' \xto{\me_2'} \sigma_3 \xfrom{\me_2} \sigma_2$ are jointly surjective.
    We show that $\concL\dash$-image of the cospan is jointly surjective which implies that the $\anyL\dash$-image is jointly surjective
    by the variable reflecting property of $\me_2'$ (\cref{lem:vrefl-prop}).

    Joint surjectivity at $\concL\dash$ is immediate from \cref{lem:alloc-epi}, every $\va : \vC \in \concL{\sigma_3} \ominus \concL{\uncon{\me_2 \comp \me_1}}$ 
    factors through $\concL{\sigma_2} \ominus \concL{\me_1}$ and so it is in the image of $\concL{\me_2}$.
\end{proof}

\begin{definition}[Category of Template Allocations]
    We consider $\allocC\tempC \subseteq \tempC$, the wide subcategory of $\tempC$ where morphisms are allocation morphisms.
    Identities are allocation morphisms and by \cref{lem:comp-alloc}, they are also closed under composition.
    Similarly, we consider $\allocC\tinstC \subseteq \tinstC$ and notice that we have a forgetful functor $\allocC\tinstC \xto{\allocC\tinstU} \allocC\tempC$.
\end{definition}

\subsubsection{Resultlets}

\emph{Resultlets} represent collections of results in the spirit of the FGLS monad.

\begin{definition}
    For a $\vA : \tinstC \to \setC$ and template $\sigma$, we define the resultlet functor $\resF\vA_\sigma$ as
    \[ \resF\vA_\sigma \eqdef \sliceC{\sigma}{\allocC\tinstU} \xto{\vU} \allocC\tinstC \inclto \tinstC \xto\vA \setC \]

    A resultlet $\vr$ over $\pair\sigma\vA$ is an element of $\typeRes\vA\sigma$
    \[ \vr \in \typeRes\vA\sigma \eqdef \Coprod{\sigma \xato{\rho'} \sigma' \in \sliceC\sigma{\allocC\tinstU}}  \resF\vA_\sigma\rho' = 
        \Coprod{\sigma \xato{\rho'} \sigma' \in \sliceC\sigma{\allocC\tinstU}} \vA(\sigma') \]
\end{definition}

\begin{remark}
For FGLS we are interested in situations where $\vA = \vX\anyL\dash \times \theapF$ for a $\vX : \worldC \to \setC$.
A resultlet over $\pair\sigma{\vX\anyL\dash \times \theapF}$ then consists of
\begin{itemize}
    \item An allocation morphism $\sigma \xato\rho \sigma'$, which by \cref{lem:codomain-alloc} 
            is determined only by $\concL\sigma \xto{\concL\rho} \concL{\sigma'}$.
        \item An $\vX$-valued return value $\vx : \vX\anyL{\sigma'}$ which by \cref{lem:codomain-alloc}
                is isomorphic to a value in $\vX(\concL{\sigma'} \oplus \varL\sigma)$.
    \item A template heap $\mvarM : \dheapF(\sigma')$, which is isomorphic to 
            a heaplet $\dheapF(\concL{\sigma'}, \concL{\sigma'} \oplus \varL\sigma)$.
\end{itemize}
\end{remark}

We define actions of $\typeRes\relax$ on morphisms in $\funC\tinstC\setC$ and $\tinstC$ but we stress that some of these actions are not functorial.
In \cref{sec:tinstC-hiding-monad} we quotient $\typeRes\relax$ to form a functor out of it.

\begin{definition}[Action on morphisms in $\funC\tinstC\setC$]
    For a natural transformation $\alpha : \vA \to \vB$, and every $\sigma \in \objC\tempC$
    \begin{align*}
        \typeRes(\alpha)_\sigma & : \typeRes\vA\sigma \to \typeRes\vB\sigma \\
        \inj_{\sigma \xato{\rho'} \sigma'}(\va) & \mapsto \inj_{\rho'}(\alpha_{\sigma'}(\va)) 
    \end{align*}
    It changes a resultlet over $\pair\sigma\vA$ to one over $\pair\sigma\vB$.
\end{definition}

\begin{definition}[Covariant action on morphisms in $\tinstC$]
For a $\me : \sigma_1 \to \sigma_2 \in \tinstC$ and a resultlet $\inj_{\sigma_1 \xato{\rho'} \sigma'} \va$
we consider the following L.I.C. in $\tempC$
\begin{center}
\begin{tikzcd}
        & {\sigma_1} 
            \ar[ld, "\tinstU\me"'] \ar[-alloctip, rd, "\rho'"] 
    \\
    {\sigma_2} 
        \ar[-alloctip, rd, "\ext^{\tinstU\me}\rho'"'] 
            && {\sigma'} 
                \ar[ld, start anchor=south west, "\ext^{\rho'}\me"] 
    \\
        & {\sigma_2 \liplus{\sigma_1} \sigma'}
\end{tikzcd}
\end{center}
Notice that by \cref{lem:lic-alloc}, $\ext^{\tinstU\me}\rho'$ is an allocation morphism.
We further equip $\ext^{\rho'}\me$ with data using \cref{rem:lic-T-data} and so define the covariant action $\typeRes\vA\me$ as
\begin{align*}
    \typeRes\vA\me & : \typeRes\vA{\sigma_1} \to \typeRes\vA{\sigma_2} \\
    \inj_{\rho'}\va & \mapsto \inj_{\ext^{\tinstU\me}{\rho'}} \vA(\ext^{\rho'}\me)(\va)
\end{align*}
\end{definition}

\begin{definition}[Contravariant action on allocation morphisms in $\tempC$]
    For every allocation morphism $\sigma_2 \xato{\rho'} \sigma_1$ in $\tempC$ define
\begin{align*}
    \hideRes\vA{\rho'} & : \typeRes\vA\sigma_1 \to \typeRes\vA\sigma_2 \\ 
    \inj_\rho \va & \mapsto \inj_{\rho \comp \rho'} \va
\end{align*}
where $\rho \comp \rho'$ is an allocation morphism itself by \cref{lem:comp-alloc}.

We call this action \emph{hiding} as it forgets some of the public locations.
\end{definition}

\begin{example}
    Consider the allocation morphism
    \[ \incl^{\va \inclto \va,\vc} \xato\rho \incl^{\va,\vb \inclto \va,\vb,\vc} \] 
    that introduces a new concrete location $\vb$,
    and a resultlet $\inj_{\rho'}\va$ where $\va : \vA(\incl^{\va,\vb \inclto \va,\vb,\vc})$.

    We can apply the covariant action $\me = \incl^{\va \inclto \va,\vc} \xto{\pair{\incl^{\va \inclto \va,\vc}}{\incl^{\va,\vc \inclto \va,\vc,\vd}}} 1_{\va,\vc}$ that instantiates $\vc$ with some data $\mvarM_\me : \dheapF(\vc, \{\va,\vc,\vd\})$

    \begin{center}
    \begin{tikzcd}
            & {\incl^{\va \inclto \va,\vc}} 
                \ar[dl, "\tinstU\me"'] \ar[dr, "\rho"] 
        \\
        {\incl^{\va,\vc \inclto \va,\vc,\vd}} 
            \ar[dr, "\ext^{\tinstU\me}\rho"']
                    && {\incl^{\va,\vb \inclto \va,\vb,\vc}} 
                        \ar[dl, "\ext^\rho\me"] 
        \\
        & {\incl^{\va,\vb,\vc \inclto \va,\vb,\vc,\vd}}
    \end{tikzcd}
    \end{center}
    to get the resultlet 
        \[ \inj_{\ext^{\tinstU\me}\rho} \vA(\ext^\rho\me)(\va) \]
            with data $\mvarM_{\ext^{\rho}{\me}}$ being $\dheapF(1, \anyL{\ext^{\tinstU\me}\rho})(\mvarM_\me)$.

    Consider the allocation morphism $\rho' = ?_\vc \xto{\pair{?_\va}{\incl^{\vc \inclto \va,\vc}}} \incl^{\va \inclto \va,\vc}$.
    We can apply the contravariant action $\rho'$ to get the resultlet $\inj_{\pair{?_{\{\va,\vb\}}}{\incl^{\vc \inclto \va,\vb,\vc}}} \va$.
\end{example}

\paragraph{Monad Structure on Resultlets}

We equip resultlets with unit and multiplication structure.

\begin{definition}[Unit of a Resultlet]
    We define unit morphisms
    \begin{align*}
        \eta^\vA_\sigma & : \vA\sigma \to \typeRes\vA\sigma \\
        \va & \mapsto \inj_{1_\sigma} \va
    \end{align*}
\end{definition}

\begin{definition}[Multiplication of a Resultlet]
    We define multiplication morphisms
    \begin{align*}
        \mu^\vA_\sigma & : \typeRes(\typeRes\vA)\sigma \to \typeRes\vA \\
        \inj_\rho \vr & \mapsto \hideRes\vA\rho\vr
    \end{align*}
\end{definition}

\subsubsection{Hiding Monad on Template Heaps}
\label{sec:tinstC-hiding-monad}

\begin{definition}[Hiding Monad on Template Heaps]
We generalise the hiding monad $\hideMon$ from \cite{kammar2017monad} to template heaps by appropriately quotienting the type of resultlets 
\[\begin{array}{r@{\spac}l@{\spac}l}
    \thideMon              & : \funC\tinstC\setC \to \funC\tinstC\setC \\
    (\thideMon\vA)(\sigma) & = \Colimit{(\resF\vA_\sigma)} \ni \eqclass{\sigma \xato{\rho'} \sigma'} \va
\end{array}\]
which satisfies the natural (monad) isomorphism for every $\vA : \tinstC \to \setC$
\[ (\thideMon\vA) \comp \inclITI\dash \hmonIso \hideMon(\vA \comp \inclITI\dash) \;:\; \instC \to \setC \]

We inject resultlets into the monad using the colimit injections $\inj^{\resF\vA_\sigma}_{\rho'} : \resF\vA_\sigma\rho' \to \thideMon\vA\sigma$
\begin{align*}
    \quot^\vA_\sigma & : \typeRes\vA\sigma \to \thideMon\vA\sigma \\
    \inj_{\rho'}\va & \mapsto \inj^{\resF\vA_\sigma}_{\rho'}(\va)
\end{align*}
\end{definition}

\begin{lemma}[Quotienting in the Hiding Monad on Template Heaps]\label{lem:thmon-quotient}
    For $\inj_{\sigma \xato{\rho_1} \sigma_1} \va_1$, and $\inj_{\sigma \xato{\rho_2} \sigma_2} \va_2$, 
    we have that $\eqclass{\rho_1}\va_1 = \eqclass{\rho_2}\va_2 \in \thideMon\vA\sigma$ iff there exist 
    \begin{tikzcd}
        & {\sigma} 
            \ar[dl, -malloctip, "\rho_1"'] \ar[dr, -alloctip, "\rho_2"]
    \\ 
    {\sigma_1} 
        \ar[dr, -alloctip, "\exists \tinstU\me_1'"'] 
            && {\sigma_2} 
                \ar[dl, -malloctip, "\exists \tinstU\me_2'"]
    \\ 
        & {\hat\sigma}
    \end{tikzcd}
    s.t. $\vA(\me_1')(\va_1) = \vA(\me_2')(\va_2)$.
\end{lemma}
\begin{proof}
    This is a direct generalisation of the proof in \cref{lem:hmon-quotient}. 
    We first show that the RHS relation is an equivalence. 
    Reflexivity and symmetry is immediate.
    For transitivity, we have
    \begin{center}
        \begin{tikzcd}
                && {\sigma_{12} \liplus{\sigma_2} \sigma_{23}}
        \\
            & {\sigma_{12}} 
                \ar[ur, -alloctip, "\ext^{\tinstU\me_{22}}\me_{23}"]
                    && {\sigma_{23}}
                        \ar[lu, -malloctip, "\ext^{\tinstU\me_{23}}\me_{22}"']
        \\
        {\sigma_1}
            \ar[ur, -malloctip, "\tinstU\me_1"]
                && {\sigma_2}
                        \ar[ul, -malloctip, "\tinstU\me_{22}"] \ar[ur, -alloctip, "\tinstU\me_{23}"']
                        && {\sigma_3}
                            \ar[ul, -malloctip, "\tinstU\me_3"']
        \\\\
                && {\sigma}
                    \ar[uull, -malloctip, "\rho_1"] \ar[uu, -alloctip, "\rho_2"] \ar[uurr, -alloctip, "\rho_3"']
        \end{tikzcd}
    \end{center}
    Notice that the top square commutes by \cref{lem:ligu-rmono-exist}
    which implies that $\inj_{\rho_1}(\va_1) \sim \inj_{\rho_3}(\va_3)$.

    To prove that the generated relation by the colimit is exactly $\sim$, one shows direction $\to$ by inducting on the list of related elements,
    and for $\from$, notice that if $\inj_{\rho_1}(\va_1) \sim \inj_{\rho_2}(\va_2)$, then
    \[ \inj_{\rho_1}\va_1 \leq \inj_{\instU\me_1 \comp \rho_1}\vA(\me_1)(\va_1) \ge \inj_{\rho_2}\va_2 \]

\end{proof}

\begin{lemma}
    $\thideMon$ together with the covariant actions on morphisms, unit and multiplication from $\typeRes\relax$ form a monad.
\end{lemma}
\begin{proof}
    \nl
    We first show it is functorial on morphisms $\me : \sigma_1 \to \sigma_2 \in \tinstC$.
    Consider two related $\inj_{\sigma_1 \xto\rho \sigma'} \va'$, $\inj_{\sigma_1 \xto{\rho'} \sigma''} \va''$. 
    By \cref{lem:thmon-quotient} we have
    \begin{tikzcd}
        & {\sigma_1} 
            \ar[dl, -malloctip, "\rho'"'] \ar[dr, -alloctip, "\rho''"] \ar[dd, "\hat\rho"]
    \\ 
    {\sigma'} 
        \ar[dr, -alloctip, "\exists \tinstU\me'"'] 
            && {\sigma''} 
                \ar[dl, -malloctip, "\exists \tinstU\me''"]
    \\ 
        & {\hat\sigma}
    \end{tikzcd}
    with $\vA(\me')(\va') = \vA(\me'')(\va'')$.

    We form the following diagram
    \begin{center}
        \begin{tikzcd}[sep=small]
            & {\hat\sigma \liplus{\sigma'} (\sigma_2 \liplus{\sigma_1} \sigma')} \ar[rr, -alloctip, "\iso'"]
                && {\hat\sigma \liplus{\sigma_1} \sigma_2}
                    && {(\sigma_2 \liplus{\sigma_1} \sigma'') \liplus{\sigma''} \hat\sigma} \ar[ll, -malloctip, "\iso''"']
        \\
        {\hat\sigma} \ar[ur, "\ext^{\tinstU\me'}\ext^{\rho'}\me"'] \ar[urrr, bend left=50, "\ext^{\hat\sigma}\me"]
                && {\sigma' \liplus{\sigma_1} \sigma_2} \ar[ul, -malloctip, "\ext^{\ext^{\rho'}\tinstU\me}\me'"]
                    && {\sigma_2 \liplus{\sigma_1} \sigma''} \ar[ur, -alloctip, "\ext^{\ext^{\rho''}\tinstU\me}\me''"']
                        && {\hat\sigma} \ar[ul, "\ext^{\tinstU\me''}\ext^{\rho''}\me"] \ar[ulll, bend right=50, "\ext^{\hat\sigma}\me"]
        \\
            & {\sigma'} \ar[ul, -malloctip, "\tinstU\me'"'] \ar[ur, "\ext^{\rho'}\me"]
                && {\sigma_2} \ar[ul, -malloctip, "\ext^{\tinstU\me}\rho'"] \ar[ur, -alloctip, "\ext^{\tinstU\me}\rho''"'] \ar[uu, -alloctip, "\ext^{\tinstU\me}\hat\rho"]
                            && {\sigma''} \ar[ur, -alloctip, "\tinstU\me''"] \ar[ul, "\ext^{\rho''}\me"']
        \\\\
                        &&& {\sigma_1} \ar[uull, -malloctip, "\rho'"'] 
                                       \ar[uu, "\tinstU\me"] 
                                       \ar[uurr, -alloctip, "\rho''"] 
                                       \ar[uuulll, -malloctip, bend left=30, "\hat\rho"] 
                                       \ar[uuurrr, -alloctip, bend right=30, "\hat\rho"']
        \end{tikzcd}
    \end{center}
    where the isomorphisms are given by the pasting property of L.I.C. in $\tempC$.
    The induced data given by \cref{rem:lic-T-data} of the morphisms also commute as they are uniquely determined by $\mvarM_\me, \mvarM_{\me'}, \mvarM_{\me''}$.

    We relate $\inj_{\ext^{\tinstU\me}\rho'}\vA(\ext^{\rho'}\me)(\va')$ with $\inj_{\ext^{\tinstU\me}\rho''}\vA(\ext^{\rho''}\me)(\va'')$ by
    \begin{center}
    \begin{tikzcd}
        & {\sigma_2} \ar[dl, -malloctip, "\ext^{\tinstU\me}\rho'"'] \ar[dr, -alloctip, "\ext^{\tinstU\me}\rho''"]
    \\
    {\sigma_2 \liplus{\sigma_1} \sigma'} \ar[dr, -alloctip, "\iso' \comp \ext^{\ext^{\rho'}\tinstU\me}\me'"']
            && {\sigma_2 \liplus{\sigma_1} \sigma''} \ar[dl, -malloctip, "\iso'' \comp \ext^{\ext^{\rho''}\tinstU\me}\me''"]
    \\ 
        & {\hat\sigma \liplus{\sigma_1} \sigma_2}
    \end{tikzcd}
    \end{center}
    which give us
    \begin{align*}
        \vA(\iso' \comp \ext^{\ext^{\rho'}\tinstU\me}\me' \comp \ext^{\rho'}\me)(\va') 
                &= \vA(\iso' \comp \ext^{\tinstU\me'}\ext^{\rho'}\tinstU\me \comp \tinstU\me')(\va') \\
                &= \vA(\ext^{\hat\sigma}\me \comp \tinstU\me'')(\va'') \\
                &= \vA(\iso'' \comp \ext^{\tinstU\me''}\ext^{\rho''}\tinstU\me \comp \tinstU\me'')(\va'') \\
                &= \vA(\iso'' \comp \ext^{\ext^{\rho''}\tinstU\me}\me'' \comp \ext^{\rho''}\me)(\va'')
    \end{align*}

    To show functoriality, for preservation of identities consider
    \begin{tikzcd}[sep=small]
        & {\sigma} \ar[dl, Rightarrow, "1"'] \ar[dr, "\rho"]
    \\
    {\sigma} \ar[dr, "\ext^{1}\rho"']
            && {\sigma'} \ar[dl, "\ext^\rho1"]
    \\
        & {\sigma \liplus\sigma \sigma'}
    \end{tikzcd}
    and so take
    \begin{tikzcd}[sep=small]
        & {\sigma} \ar[dl, "\ext^{1}\rho"'] \ar[dr, "\rho"]
    \\
    {\sigma \liplus\sigma \sigma'} \ar[dr, Rightarrow, "1"']
            && {\sigma'} \ar[dl, "\ext^\rho1"]
    \\
        & {\sigma \liplus\sigma \sigma'}
    \end{tikzcd}
    and for composition one considers pasting law for L.I.C. 

    That multiplication is quotient respecting, consider
    \begin{tikzcd}
        & {\sigma} \ar[dl, "\rho_1'"'] \ar[dr, "\rho_2'"]
    \\
    {\sigma_1'} \ar[dr, "\tinstU\me_1"']
            && {\sigma_2'} \ar[dl, "\tinstU\me_2"]
    \\
        & {\hat\sigma}
    \end{tikzcd}
    for some $\vq_1 = \quot_{\sigma_1' \xto{\rho_1''} \sigma_1''} \va_1''$, and  
    $\vq_2 = \quot_{\sigma_2' \xto{\rho_2''} \sigma_2''} \va_2''$.
    with
    \[ \thideMon(\vA)(\me_1)(\vq_1) = \thideMon(\vA)(\me_2)(\vq_2) \]

    Unpacking this gives us a diagram
    \begin{center}
        \begin{tikzcd}
                && {\sigma} \ar[dl, "\rho_1'"'] \ar[dr, "\rho_2'"]
        \\
            & {\sigma_1'} \ar[dr, "\tinstU\me_1"] \ar[dl, "\rho_1''"']
                        && {\sigma_2'} \ar[dl, "\tinstU\me_2"'] \ar[dr, "\rho_2''"]
        \\
        {\sigma_1''} \ar[dr, "\ext^{\rho_1''}\me_1"']
                && {\hat\sigma} \ar[dl, "\ext^{\tinstU\me_1}\rho_1''"'] \ar[dr, "\ext^{\tinstU\me_2}\rho_2''"]
                        && {\sigma_2''} \ar[dl, "\ext^{\rho_2''}\me_2"]
        \\ 
            & {\sigma_1'' \liplus{\sigma_1'} \hat\sigma} \ar[dr, "\tinstU\check\me_1"']
                    && {\sigma_2'' \liplus{\sigma_2'} \hat\sigma} \ar[dl, "\tinstU\check\me_2"]
        \\
                && {\check\sigma}
        \end{tikzcd}
    \end{center}
with $\vA(\tinstU\check\me_1 \comp \ext^{\rho_1'}\me_1)(\va_1'') = \vA(\tinstU\check\me_2 \comp \ext^{\rho_2''}\me_2)(\va_2'')$. 

    Which is exactly what we need to show that 
    \[ \hideRes\vA{\rho_1'}\vq_1 = \hideRes\vA{\rho_2'}\vq_2 \]

    That the action on morphisms in $\funC\tinstC\setC$ is quotient respecting and natural is a simple exercise
    and the same goes for naturality of unit and multiplication and that they satisfy the monad laws.
\end{proof}

\begin{lemma}\label{lem:thidemon-hide-ext-com}
    Given a L.I.G.U. diagram
    \begin{tikzcd}[sep=small]
        & {\sigma} 
            \ar[dl, "\me_1"'] \ar[dr, -alloctip, "\me_2"] 
    \\ 
    {\sigma_1} 
        \ar[dr, -alloctip, "\ext^{\me_1}\me_2"']
            && {\sigma_2}
                \ar[dl, "\ext^{\me_2}\me_1"] 
    \\ 
        & {\sigma_1 \liplus\sigma \sigma_2}
    \end{tikzcd} and a $\vr \in \thideMon\vA(\sigma_2)$, we have
    \[ \hideRes\vA{\ext^{\me_1}\me_2}(\thideMon(\ext^{\me_2}\me_1)(\vr)) = \thideMon(\me_1)(\hideRes\vA{\me_2}(\vr)) \] 
\end{lemma}
\begin{proof}
    Let $\vr = \eqclass{\sigma_2 \xato\rho \hat\sigma} \va$, for some $\va : \vA(\hat\sigma)$.
    The LHS of the equation gives
    \[ 
        \eqclass{\sigma_1 \xato{\tinstU(\ext^\me_1\me_2)} \sigma_1 \liplus\sigma \sigma_2 
                          \xato{\ext^{\tinstU(\ext^{\me_2}\me_1)}\rho)} (\sigma_1 \liplus\sigma \sigma_2) \liplus{\sigma_2} \hat\sigma} 
                            \vA(\ext^\rho(\ext^{\me_2}\me_1)(\va))
    \]
    using the diagram
    \begin{center}
        \begin{tikzcd}[sep=small]
            & {\sigma_2} \ar[dr, -alloctip, "\rho"] \ar[dl, "\ext^{\me_2}\me_1"']
        \\
        {\sigma_1 \liplus\sigma \sigma_2} \ar[dr, -alloctip, "\ext^{\tinstU(\ext^{\me_2}\me_1)}\rho"']
                && {\hat\sigma} \ar[dl, "\ext^\rho(\ext^\me_2\me_1)"]
        \\
            & {(\sigma_1 \liplus\sigma \sigma_2) \liplus{\sigma_2} \hat\sigma}
        \end{tikzcd}
    \end{center}

    The RHS of the equation gives
    \[ 
        \eqclass{\sigma_1 \xato{\ext^{\tinstU\me_1}(\rho \comp \tinstU\me_2)} \sigma_1 \liplus\sigma \hat\sigma} \vA(\ext^{(\rho \comp \tinstU\me_2)}\me_1)(\va)
    \]
    using the diagram
    \begin{center}
        \begin{tikzcd}[sep=small]
                && {\sigma} \ar[dr, -alloctip, "\tinstU\me_2"] \ar[ddll, "\me_1"']
        \\
                    &&& {\sigma_2} \ar[dr, -alloctip, "\rho"]
        \\
        {\sigma_1} \ar[ddrr, -alloctip, "\ext^{\tinstU\me_1}(\rho \comp \tinstU\me_2)"']
                        &&&& {\hat\sigma} \ar[ddll, "\ext^{(\rho \comp \tinstU\me_2)}\me_1"]
        \\
        \\
                && {\sigma_1 \liplus\sigma \hat\sigma}
        \end{tikzcd}
    \end{center}

    We relate the two resultlets by
    \begin{center}
        \begin{tikzcd}
            & {\sigma} \ar[dl, "\me_1"'] \ar[dr, -alloctip, "\tinstU\me_2"]
        \\
        {\sigma_1} \ar[dr, -alloctip, "\tinstU(\ext^{\me_1}\me_2)"'] \ar[drr, -alloctip, "\ext^{\tinstU\me_1}(\rho \comp \tinstU\me_2)" near start]
                    && {\sigma_2} \ar[dr, -alloctip, end anchor=north east, "\rho"] \ar[dl, crossing over, "\ext^{\me_2}\me_1"']
        \\
                & {\sigma_1 \liplus\sigma \sigma_2} \ar[dr, -alloctip, "\ext^{\tinstU(\ext^{\me_2}\me_1)}\rho"']
                    & {\sigma_1 \liplus\sigma \hat\sigma} \ar[d, dashed, "\iso"', "\vh"]
                        & {\hat\sigma} \ar[l, "\ext^{(\rho \comp \tinstU\me_2)}\me_1"'] \ar[dl, "\ext^\rho(\ext^{\me_2}\me_1)"]
        \\
                && {(\sigma_1 \liplus\sigma \sigma_2) \liplus{\sigma_2} \hat\sigma}
        \end{tikzcd}
    \end{center}
    where by L.I.C. pasting (\cref{lem:lic-pasting}) the two outer L.I.C. squares are isomorphic,
    witnessed by the iso $\vh$.
    Further, $\vh$ is an allocation morphism by \cref{lem:alloc-cancl}.
    And so, by \cref{lem:thmon-quotient}, the resultlets are related via $\vh$ and the identity $1_{(\sigma_1 \liplus\sigma \sigma_2) \liplus{\sigma_2} \hat\sigma}$.
    (notice that the data of $\ext^\rho(\ext^{\me_2}\me_1)$ and $\ext^{(\rho \comp \tinstU\me_2)}\me_1$ are completely determined by $\mvarM_{\me_1}$ and so they match).
\end{proof}

\subsubsection{Evaluation and Realisation}

An \emph{evaluation} is a collection of resultlets. 

\begin{definition}[Evaluation]
    Let $\Sigma$ be a finite family of templates and $\vA : \tinstC \to \setC$.
    An $\vA$-valued \emph{evaluation} $\ev$ over $\Sigma$ is an indexed family of resultlets $\ev \in \Prod{\sigma_\vi \in \Sigma} \typeRes\vA\sigma_\vi$.
\end{definition}

Let $\vA : \tinstC \to \setC$, total cover $\cov$ over $\pair\vw\Sigma$, $\vA$-valued evaluation $\ev$ over $\Sigma$, 
and let $\phi$ be
\begin{center}
    \begin{tikzcd}
        {\homF\worldC\vw{\instU(\dash)} \times \heapF} \ar[r, "\phi"]
            & {\hideMon{(\vA \comp \inclITI\dash)}} \ar[d, "\hmonIso"] 
        \\
        {\stencilF\vw\inclITI\dash} \ar[r, dashed] \ar[u, "\pair{\anyL\dash}{\mvarM_\dash}", "\iso"']
            & {\thideMon\vA\inclITI\dash}
    \end{tikzcd}
\end{center}
which induces a morphism from $\stencilF\vw\inclITI\dash \dashto \thideMon\vA\inclITI\dash$.

\begin{definition}[Realisation]
We say that $\cov$ and $\ev$ realise $\phi$, $\realise\cov\ev\phi$ if
for every concrete stencil $\me :\, ?_\vw \to 1_{\hat\vw}$, we have
\[ \hmonIso(\phi\,\pair{\anyL\me}{\mvarM_\me}) = \thideMon(\vA)(\tilde\me)(\proj_{\indF{\tilde\me}(\star)} \ev) \]  
where the morphism $\tilde\me : \sigma_{\indF{\tilde\me}(\star)} \to 1_{\hat\vw}$ is given by totality.
\end{definition}
